
% #### 2. Nástroj 1: Curl Dashboard (`curl.se/dashboard.html`)
% * **Charakteristika:** Curl je masivní, desítky let starý projekt. Jejich dashboard je ultimátní ukázka "broad-spectrum" monitoringu.
% * **Co to umí (Zaměření na kvantitu):** Sleduje úplně všechno od počtu commitů za měsíc, přes statistiky hlášení zranitelností (CVE), pokrytí kódu testy (code coverage), až po počet aktivních autorů.
% * **Forma:** Je to obří agregátor grafů.
% * **Zhodnocení:** * *Plusy:* Neuvěřitelně komplexní historická data, skvělé pro manažerský přehled o "zdraví" projektu v dlouhodobém horizontu.
%     * *Mínusy:* Je to spíše statický report (statistiky) než interaktivní nástroj pro operativní řešení záseků (např. nepomůže ti zjistit, který PR právě teď blokuje CI pipelinu).

\section{Curl Dashboard}
\label{sec:curl_dashboard}

Curl Dashboard\footnote{\url{https://curl.se/dashboard.html}} je statický nástroj pro zobrazení různých statistik o~projektu curl za celou dobu jeho existence. Jeho primárním účelem není operativní řízení vývoje, nýbrž slouží jako zdroj informací z~dlouhodobého hlediska pro sledování zdraví projektu a~jeho trendů.

\subsection{Obsah a~funkcionalita}
\label{subsec:curl_dashboard:content}

Jedná se o~statickou stránku zobrazující denně generované grafy a~statistiky o~projektu curl. Tyto grafy jsou ve formátu obrázků a~jsou generovány pomocí skriptů analyzujících data z~git repozitáře a~\API GitHubu projektu curl.

Na rozdíl od GitHub Insights \chapterref{sec:github_stats} zde není interaktivní přístup k~datům ani možnost výběru zobrazovaného časového rozsahu. Kromě dat dostupných přímo na GitHubu jsou zde zobrazeny i~další informace, jako například statistiky bug bounty programu, počty chyb a~bezpečnostních zranitelností, distribuce commitů podle časových pásem nebo poměr počtu úkolů (TODO) k~počtu známých chyb.

\subsection{Zhodnocení}
\label{subsec:curl_dashboard:evaluation}

Curl Dashboard poskytuje velmi komplexní a~detailní pohled na dlouhodobý stav a~zdraví projektu curl. Hlavní výhodou je velké množství grafů a~statistik pokrývajících širokou škálu aspektů projektu. Samotné grafy jsou jednoduché, přehledné a~srozumitelné.

Nevýhodou tohoto nástroje je absence interaktivity. Velké množství grafů na jedné stránce může ztěžovat orientaci a~nalezení požadovaných informací. Protože se jedná o~čistě statickou stránku, není možné získat detailní informace o~konkrétním dni nebo měsíci a při výskytu výrazné odchylky v~grafu nelze zjistit, co přesně se stalo.
